\documentclass[11pt]{amsart}
\usepackage[T1]{fontenc}
\usepackage{lmodern,amsmath,amssymb,amsthm}
\usepackage[all]{xy}
\usepackage[hidelinks]{hyperref}
\emergencystretch=2em
\ifdefined\pdfinfoomitdate\pdfinfoomitdate=1\fi
\ifdefined\pdftrailerid\pdftrailerid{}\fi
\newtheorem{proposition}{Proposition}[section]
\newtheorem{lemma}[proposition]{Lemma}
\newtheorem{definition}[proposition]{Definition}
\theoremstyle{remark}
\newtheorem{remark}[proposition]{Remark}
\newcommand{\PSh}{\operatorname{PSh}}
\newcommand{\El}{\operatorname{El}}
\newcommand{\Fun}{\operatorname{Fun}}
\newcommand{\Hom}{\operatorname{Hom}}
\newcommand{\Cat}{\mathbf{Cat}}
\newcommand{\Sets}{\mathbf{Sets}}
\newcommand{\id}{\operatorname{id}}
\newcommand{\PSM}[1]{\par\smallskip{\small\raggedright\noindent Source-map identifiers: \texttt{#1}.\par}\smallskip}
\title{Categories of Elements, Relative Nerves, and a Counit Criterion}
\author{Unofficial Stacks Project AI Drafts}
\date{September 2026}
\begin{document}
\maketitle
\begin{abstract}
An independently readable module adapted from a source-mapped draft of
Grothendieck's \emph{Pursuing Stacks}, sections 19, 26, 28--30, and 38.
We give explicit functors and proofs for a realization--nerve adjunction,
the category-of-elements adjunction, its comma-category description, and
a conditional criterion for an equivalence after localization.
This is AI-written draft mathematics, not official Stacks text or
proof-assistant-verified mathematics. Existing results are identified as
dependencies, not claimed as new theorems.
\end{abstract}
\tableofcontents

\section{Conventions and existing foundations}
\label{psm-conventions}

Fix a universe. Small categories and sets below belong to it; functor
categories may be formed in a larger universe. The notation $\Fun(C,D)$
denotes the category of functors and natural transformations when used as
a category, and its set of objects when it occurs in a set-valued
presheaf. Thus morphisms in the ordinary category $\Cat$ are functors,
not natural-isomorphism classes of functors.

For a small category $A$, write $\PSh(A)=\Fun(A^{\mathrm{op}},\Sets)$.
For $F\in\PSh(A)$ the category $\El_A(F)$ has objects $(a,x)$,
where $x\in F(a)$, and arrows
\[
 (a,x)\xrightarrow{u}(b,y)\quad\text{with}\quad
 u:a\longrightarrow b,\qquad F(u)(y)=x.
\]
It is the category fibred in sets already constructed in Stacks,
\nolinkurl{categories.tex}, \nolinkurl{example-presheaf}.
We use it rather than introduce a competing definition.
For a representable $h_a=\Hom_A(-,a)$, it is canonically $A/a$.

We also use the existing Yoneda lemma and the presentation
\begin{equation}\label{psm-density}
 F\simeq\mathop{\mathrm{colim}}_{(a,x)\in\El_A(F)}h_a.
\end{equation}
The proof is the presheaf calculation in the following Stacks source:
\begin{center}\small
\nolinkurl{sites.tex}: \nolinkurl{lemma-colimit-representable}.
\end{center}
At $b$, the class of $(a,x,u:b\to a)$ maps to $F(u)(x)$,
and $y\in F(b)$ is represented by $(b,y,\id_b)$.
Taking the indiscrete topology gives the presheaf version directly.
These foundations are existing coverage, not contributions of this module.

\section{Presheaves over a presheaf}
\label{psm-slice-section}

\begin{proposition}\label{psm-slice-equivalence}
For $F\in\PSh(A)$ there is a natural equivalence
\[
 \PSh(A)/F\simeq\PSh(\El_A(F)).
\]
It sends $p:E\to F$ to the presheaf with value $p_a^{-1}(x)$ at $(a,x)$.
\end{proposition}
\PSM{PSM-PROP-0008}
\begin{proof}
If $u:(a,x)\to(b,y)$, naturality of $p$ says that $E(u)$ carries
$p_b^{-1}(y)$ into $p_a^{-1}(x)$. This defines the asserted presheaf.
Conversely, for $G\in\PSh(\El_A(F))$, set
\[
 E_G(a)=\coprod_{x\in F(a)}G(a,x).
\]
For $u:a\to b$ send $(y,s)$ to
$(F(u)(y),G(\widetilde u)(s))$, where
$\widetilde u:(a,F(u)(y))\to(b,y)$ is the arrow over $u$.
Functoriality follows from that of $F$ and $G$.
Projection onto $x$ defines $E_G\to F$.
Maps over $F$ induce maps on fibres, and a natural transformation between
two such $G$ induces a map on the displayed disjoint unions. These rules
are functorial. The fibre of $E_G(a)\to F(a)$ at $x$ is canonically
$G(a,x)$, while the disjoint union of the fibres of $p_a$ is $E(a)$.
The resulting two isomorphisms commute with all restriction maps and
with maps between presheaves. They give the equivalence.
\end{proof}

\begin{remark}
This is an explicit presheaf formulation of localization of a topos,
not a claim of a new localization theorem. Compare Stacks,
\nolinkurl{sites.tex}, \nolinkurl{lemma-localize-topos} and
\nolinkurl{remark-localize-presheaves}. The additional value here is the
fixed category-of-elements model and the two inverse formulas used below.
\end{remark}

\section{Relative nerves and free cocompletion}
\label{psm-realization-section}

Let $D:A\to M$ be a functor, where $M$ is locally small and has all
small colimits. Define
\[
 N_D(X)(a)=\Hom_M(D(a),X),\qquad
 D_!(F)=\mathop{\mathrm{colim}}_{(a,x)\in\El_A(F)}D(a).
\]
The restriction in $N_D(X)$ is precomposition with $D(u)$.
A map of presheaves induces a map of the indexing categories over $A$,
and hence makes $D_!$ a functor.

\begin{proposition}\label{psm-realization-adjunction}
There is a natural adjunction $D_!\dashv N_D$. The functor $D_!$
preserves small colimits and has a canonical isomorphism $D_!h_a\simeq D(a)$.
Restriction along Yoneda induces an equivalence of categories
\[
 \Fun_{\mathrm{colim}}(\PSh(A),M)\ \simeq\ \Fun(A,M),
\]
where the left-hand category has small-colimit-preserving functors as
objects and all natural transformations as morphisms.
In particular, an extension of a specified $D$, together with its specified
identification on representables, is unique up to a unique compatible
natural isomorphism.
\end{proposition}
\PSM{PSM-PROP-0004; PSM-PROP-0031}
\begin{proof}
A map $D_!(F)\to X$ is a family of maps $D(a)\to X$, indexed by
$x\in F(a)$, such that for $u:a\to b$ and $y\in F(b)$ the map
indexed by $F(u)(y)$ is the map indexed by $y$ composed with $D(u)$.
That is exactly a natural transformation $F\to N_D(X)$. The
construction is inverse in both directions and natural in both variables.
Thus $D_!$ is a left adjoint and preserves small colimits.
Since $\El_A(h_a)=A/a$ has terminal object $\id_a$, its defining colimit
is canonically $D(a)$.

For the last assertion, a natural transformation $D\to D'$ induces,
by taking the displayed colimits, a transformation $D_!\to D'_!$.
For any colimit-preserving $L$, equation~\eqref{psm-density} gives a
natural isomorphism
\[
 L(F)\simeq\mathop{\mathrm{colim}}_{(a,x)\in\El_A(F)}L(h_a).
\]
Consequently every transformation on the representables extends uniquely
to all presheaves. This proves full faithfulness of restriction. The
construction $D\mapsto D_!$ proves essential surjectivity. Fixing the
identification on representables then gives the stated uniqueness.
\end{proof}

\begin{remark}
The compatibility clause in the uniqueness assertion is essential. An
unspecified isomorphism $L\circ y\simeq D$ does not eliminate natural
automorphisms of $D$. The functor-category equivalence states the precise
universal property. Specializing to $M=\Cat$ supplies the generalized
realization--nerve adjunction of source section 38; no weak-equivalence
property follows merely from this adjunction.
\end{remark}

\section{The right adjoint to categories of elements}
\label{psm-elements-section}

For a small category $C$ put
\[
 J_A(C)(a)=\Fun(A/a,C).
\]
For $u:a\to b$, restriction is precomposition with
$u_*:A/a\to A/b$, which sends $v$ to $uv$.

\begin{proposition}\label{psm-elements-adjunction}
The functors $\El_A:\PSh(A)\to\Cat$ and $J_A:\Cat\to\PSh(A)$
form an adjoint pair $\El_A\dashv J_A$.
\end{proposition}
\PSM{PSM-DEF-0030; PSM-PROP-0011}
\begin{proof}
Given $H:\El_A(F)\to C$ and $x\in F(a)$, associate the functor
\[
 H_x:A/a\longrightarrow C,\qquad
 (v:b\to a)\longmapsto H(b,F(v)(x)).
\]
On a triangle in $A/a$ use the induced arrow in $\El_A(F)$.
The identity $H_{F(u)(y)}=H_y\circ u_*$ makes $x\mapsto H_x$
a natural transformation $F\to J_A(C)$.

Conversely, given $\theta:F\to J_A(C)$, set
$H(a,x)=\theta_a(x)(\id_a)$. An arrow
$u:(a,x)\to(b,y)$ satisfies $x=F(u)(y)$. Thus
$\theta_a(x)=\theta_b(y)\circ u_*$, and the arrow
$u\to\id_b$ in $A/b$ gives
\[
 H(a,x)=\theta_b(y)(u)\longrightarrow\theta_b(y)(\id_b)=H(b,y).
\]
Identities and compositions follow by evaluating the corresponding
triangles in a slice. The first construction followed by the second
recovers $H$ on objects and arrows. In the other order, at an object
$v:b\to a$ naturality identifies the recovered value with
$\theta_a(x)(v)$, and does the same for slice arrows. Hence it recovers
$\theta$ exactly. This is a natural bijection of hom sets.
\end{proof}

Write $E_A(C)=\El_A(J_A(C))$. The counit is the evaluation functor
\begin{equation}\label{psm-counit}
 \epsilon_C:E_A(C)\longrightarrow C,\qquad
 (a,T:A/a\to C)\longmapsto T(\id_a).
\end{equation}
An arrow over $u:a\to b$ satisfies $T=S\circ u_*$; its image is
$S(u\to\id_b)$. This makes evaluation natural in $C$.

\begin{lemma}\label{psm-counit-comma}
For $c\in C$ there is a canonical isomorphism of categories
\[
 (\epsilon_C\mathbin\downarrow c)\ \cong\ E_A(C/c).
\]
\end{lemma}
\PSM{PSM-LEM-0018}
\begin{proof}
An object on the left is $(a,T,q)$ with $T:A/a\to C$ and
$q:T(\id_a)\to c$. It defines $\widetilde T:A/a\to C/c$ by
\[
 v:b\to a\ \longmapsto\
 \bigl(T(v)\xrightarrow{T(v\to\id_a)}T(\id_a)\xrightarrow{q}c\bigr).
\]
Conversely, compose a functor $A/a\to C/c$ with the forgetful functor
and retain its structural arrow at $\id_a$. Naturality of structural
arrows in the slice shows these operations are inverse on objects.
An arrow on the left lies over $u:a\to b$, satisfies $T=S u_*$, and
satisfies $q=r\,S(u\to\id_b)$. These equalities are precisely
$\widetilde T=\widetilde S u_*$, the condition for an arrow on the right.
Thus the operations are inverse on morphisms too and preserve composition.
\end{proof}

\section{A conditional counit-test criterion}
\label{psm-counit-criterion-section}

Let $W$ be a class of functors between small categories. Say that $C$ is
$W$-aspheric if $C\to\mathbf 1$ belongs to $W$. Say that $f:C\to D$
is right $W$-aspheric if $(f\downarrow d)$ is $W$-aspheric for every $d$.
We assume explicitly that:
\begin{enumerate}
\item $W$ contains all isomorphisms and satisfies two out of three;
\item a category with a final object is $W$-aspheric;
\item every right $W$-aspheric functor belongs to $W$.
\end{enumerate}
These are hypotheses on $W$, not claims about an arbitrary proposed
notion of weak equivalence. The last condition is the relevant
Theorem-A-type input; it is not proved by the criterion below.

\begin{proposition}\label{psm-counit-test}
For a small category $A$, the following conditions are equivalent:
\begin{enumerate}
\item $\epsilon_C\in W$ for every small $C$;
\item $\epsilon_C$ is right $W$-aspheric for every small $C$;
\item $E_A(B)$ is $W$-aspheric whenever $B$ has a final object.
\end{enumerate}
\end{proposition}
\PSM{PSM-DEF-0032; PSM-DEF-0033; PSM-DEF-0034; PSM-PROP-0013}
\begin{proof}
If (1) holds and $B$ has a final object, both $\epsilon_B$ and
$B\to\mathbf 1$ belong to $W$, so their composite does. This is (3).
Assume (3). For $c\in C$, the category $C/c$ has final object
$\id_c$. Lemma~\ref{psm-counit-comma} identifies
$(\epsilon_C\downarrow c)$ with $E_A(C/c)$, which is $W$-aspheric
by (3). The isomorphism does not change membership of the terminal
functor in $W$, by (1) of the hypotheses and two out of three.
Thus (2) holds. The final hypothesis on $W$ gives (2)$\Rightarrow$(1).
\end{proof}

\noindent The proof mechanism, with its hypotheses visible, is
\[
 \xymatrix@C=30pt{
 E_A(C/c)\ar[r]^-{\sim}\ar[dr] &
 (\epsilon_C\downarrow c)\ar[d]\\
 &\mathbf 1.
 }
\]
The top arrow is the explicit isomorphism of
Lemma~\ref{psm-counit-comma}. A final object in $C/c$ and
condition (3) make the diagonal terminal map weak, hence also the vertical map;
the Theorem-A-type hypothesis then makes $\epsilon_C$ weak.

\section{Passage through localization}
\label{psm-localization-section}

\begin{lemma}\label{psm-localization-counit}
Let $L:M\rightleftarrows N:R$ be an adjunction and let $W_N$ contain
identities and satisfy two out of three. Put $W_M=L^{-1}(W_N)$.
If every component of the counit $\epsilon:LR\to\id_N$ belongs to
$W_N$, then $R(W_N)\subset W_M$, and every component of the unit
$\eta:\id_M\to RL$ belongs to $W_M$. If the ordinary categorical
localizations at these classes exist, $L$ and $R$ induce mutually
quasi-inverse equivalences between them.
\end{lemma}
\PSM{PSM-LEM-0016}
\begin{proof}
For $f:Y\to Y'$ in $W_N$, the square
\[
 \xymatrix{
 LRY\ar[r]^{LR(f)}\ar[d]_{\epsilon_Y}&LRY'\ar[d]^{\epsilon_{Y'}}\\
 Y\ar[r]_f&Y'
 }
\]
commutes. Its bottom and vertical arrows belong to $W_N$; applying two
out of three twice gives $LR(f)\in W_N$, hence $R(f)\in W_M$.
The triangle identity
$\epsilon_{LX}\circ L(\eta_X)=\id_{LX}$ gives $L(\eta_X)\in W_N$,
so $\eta_X\in W_M$. The universal properties of localization now
give induced functors. The unit and counit descend to natural
transformations: naturality for an inverted arrow follows from naturality
for that arrow by multiplying by its inverse. All their components become
isomorphisms. The triangle identities remain identities after localization,
so the induced functors are quasi-inverse adjoint equivalences.
\end{proof}

\begin{proposition}\label{psm-test-localization}
Assume the hypotheses and equivalent conditions of
Proposition~\ref{psm-counit-test}. On $\PSh(A)$ set
$W_A=\{u:\El_A(u)\in W\}$. If the two localizations exist, then
\[
 \PSh(A)[W_A^{-1}]\ \simeq\ \Cat[W^{-1}],
\]
with functors induced by $\El_A$ and $J_A$.
\end{proposition}
\begin{proof}
Apply Lemma~\ref{psm-localization-counit} to the adjunction of
Proposition~\ref{psm-elements-adjunction}. Its counit is precisely
\eqref{psm-counit}, which lies in $W$ by the assumed criterion.
\end{proof}

\begin{remark}
This last conclusion is conditional. It neither constructs the
localizations nor asserts that any particular $A$ satisfies the criterion.
No model structure, comparison with spaces, or arbitrary-topos theorem is
silently imported. The term \emph{counit-test} specifies this criterion;
it does not identify all provisional meanings of ``test category'' in
the historical source.
\end{remark}

\section*{Source and reuse notes}
The historical source is Alexander Grothendieck, \emph{Pursuing Stacks},
\emph{First Episode: The Modelizing Story}, source edition
\href{https://arxiv.org/abs/2111.01000}{arXiv:2111.01000v2}.
The source-map identifiers printed here are local to the received draft,
not official Stacks tags. Exact source spans, received-file hashes,
comparison decisions, and the scope of the checks accompany this module
in \texttt{integration.json} and \texttt{REVIEW.md}.

The Stacks comparison baseline is
\href{https://github.com/stacks/stacks-project/tree/a04446e57ec1fbc252a871afcec7752fb2807b14}{\texttt{a04446e57ec1}};
existing FGA nerve coverage is recorded separately. This is additional
exposition, not a claim of novelty or upstream errata. This complete
LaTeX file uses only standard packages. Its source archive supplies the
build tools and retained licensing and attribution notices.
\end{document}
